% !TEX program = xelatex
% ============================================================
%  خطة تحسين تصميم الفون — تجربة الجوال عبر كل الصفحات
%  الاستخدام على Overleaf:
%    1) New Project ← Blank Project
%    2) الصق المحتوى في main.tex
%    3) Menu ← Compiler ← اختر: XeLaTeX   (مهم للعربية)
%    4) Recompile ← نزّل PDF
% ============================================================
\documentclass[11pt]{article}

\usepackage{polyglossia}
\setmainlanguage{arabic}
\setotherlanguage{english}
\newfontfamily\arabicfont[Script=Arabic,Scale=1.0]{Amiri}
\newfontfamily\arabicfontsf[Script=Arabic,Scale=1.0]{Amiri}
\setmonofont{DejaVu Sans Mono}[Scale=0.8]

\usepackage{geometry}
\geometry{a4paper, margin=2.2cm}

\usepackage{xcolor}
\usepackage{graphicx}
\usepackage{booktabs}
\usepackage{longtable}
\usepackage{array}
\usepackage{enumitem}
\usepackage{titlesec}
\usepackage{fancyhdr}
\usepackage{hyperref}
\usepackage{tcolorbox}
\tcbuselibrary{skins,breakable}
\usepackage{fontawesome5}

% ── الألوان (هوية يَسِّر) ──
\definecolor{gold}{HTML}{B8933B}
\definecolor{teal}{HTML}{0F3D3E}
\definecolor{darkbg}{HTML}{101414}
\definecolor{okgreen}{HTML}{27AE60}
\definecolor{warnred}{HTML}{C0392B}
\definecolor{softgray}{HTML}{F4F4F4}
\definecolor{midgray}{HTML}{7A7A7A}
\definecolor{codebg}{HTML}{1A2020}

\pagestyle{fancy}
\fancyhf{}
\renewcommand{\headrulewidth}{0.4pt}
\renewcommand{\footrulewidth}{0.4pt}
\fancyhead[R]{\small\color{gold}خطة تحسين تصميم الفون — منصة يَسِّر}
\fancyhead[L]{\small\color{midgray}Mobile UX Enhancement}
\fancyfoot[C]{\small\thepage}

\titleformat{\section}{\Large\bfseries\color{gold}}{\thesection}{0.6em}{}
\titleformat{\subsection}{\large\bfseries\color{darkbg}}{\thesubsection}{0.5em}{}
\titleformat{\subsubsection}{\normalsize\bfseries\color{midgray}}{\thesubsubsection}{0.4em}{}

\setlist[itemize]{rightmargin=0pt, itemsep=2pt, topsep=3pt}
\setlist[enumerate]{rightmargin=0pt, itemsep=2pt, topsep=3pt}

\hypersetup{
  colorlinks=true, linkcolor=gold, urlcolor=gold,
  pdftitle={خطة تحسين تصميم الفون}, pdfauthor={منصة يَسِّر}
}

% ── صناديق مخصّصة ──
\newtcolorbox{goalbox}[1][]{
  colback=gold!8, colframe=gold, boxrule=0.8pt,
  arc=3pt, left=8pt, right=8pt, top=6pt, bottom=6pt,
  breakable, #1
}
\newtcolorbox{notebox}[1][]{
  colback=softgray, colframe=midgray, boxrule=0.6pt,
  arc=3pt, left=8pt, right=8pt, top=6pt, bottom=6pt,
  breakable, #1
}
\newtcolorbox{warnbox}[1][]{
  colback=warnred!6, colframe=warnred, boxrule=0.8pt,
  arc=3pt, left=8pt, right=8pt, top=6pt, bottom=6pt,
  breakable, #1
}
\newtcolorbox{codebox}[1][]{
  colback=codebg, colframe=codebg, coltext=white,
  arc=3pt, left=8pt, right=8pt, top=6pt, bottom=6pt,
  breakable, #1
}

\begin{document}

% ============================================================
\begin{center}
  {\Huge\bfseries\color{gold} خطة تحسين تصميم الفون}\\[4pt]
  {\large\color{darkbg} تجربة جوال موحّدة عبر كل صفحات المنصة}\\[6pt]
  {\small\color{midgray} منصة يَسِّر الفقهية — \today}
\end{center}

\vspace{4pt}
\hrule
\vspace{10pt}

% ============================================================
\section{الهدف}

\begin{goalbox}
توحيد تجربة الاستخدام على الجوال عبر كل الصفحات على أساس \textbf{النهج الحديث} المُتَّبع جزئيًا في
صفحة الـ Hero (\texttt{clamp()}, \texttt{dvh}, \texttt{safe-area-inset}، مناطق لمس مريحة)،
ونقل هذه المعايير إلى الصفحات والمكوّنات المتبقية لتصبح المنصة متّسقة وسريعة ومريحة على الهاتف.
\end{goalbox}

\subsection*{نطاق العمل}
\begin{itemize}
  \item اعتماد \textbf{Mobile-First} فعليًا بدل استثناءات لاحقة في نهاية الملفات.
  \item توحيد \textbf{مناطق اللمس}، \textbf{المسافات الآمنة}، \textbf{أحجام النص}، \textbf{شريط التنقل}.
  \item تحسين الأداء والإحساس الأصلي عند التثبيت كتطبيق (PWA).
\end{itemize}

% ============================================================
\section{الوضع الحالي}

\begin{notebox}
\textbf{نقاط القوة الموجودة:}
\begin{itemize}
  \item صفحة الـ Hero متفوّقة: \texttt{dvh}، \texttt{safe-area}، \texttt{clamp}، \texttt{prefers-reduced-motion}.
  \item \texttt{viewport-fit=cover} مضبوط في \texttt{index.html}، وشريط التنقل السفلي يدعم الـ Safe Area.
  \item خطوط محلية (\texttt{Tajawal}, \texttt{Amiri}) بلا اعتماد على شبكات خارجية.
\end{itemize}

\vspace{4pt}
\textbf{نقاط الضعف:}
\begin{itemize}
  \item أنماط كبيرة مسطّحة: \texttt{index.css} ($\sim$1442 سطرًا)، \texttt{exterior.css} ($\sim$866 سطرًا)
        مع \texttt{!important} متكرّر يُصعّب الضبط.
  \item نصوص بأحجام ثابتة (\texttt{fontSize}) عبر \texttt{inline styles} في مكوّنات عدّة، لا تستجيب لحجم الشاشة.
  \item مناطق لمس قد تكون أصغر من الحد الموصى به (44$\times$44px) في بعض الأزرار.
  \item أنماط استجابة غير متّسقة: بعض المكوّنات لها \texttt{@media} وبعضها لا.
  \item ألوان الظلّ ثابتة (\texttt{rgba(0,0,0,...)}) تظلّ ثقيلة/باهتة في وضع الجوال الداكن.
\end{itemize}
\end{notebox}

% ============================================================
\section{المرحلة الأولى: نظام رموز التصميم المتجاوب}

إنشاء متغيّرات موحّدة في \texttt{:root} داخل \texttt{index.css} تُستخدم في كل المكوّنات:

\begin{codebox}
{\ttfamily\small\raggedleft
:root \{\\
\quad /* مناطق اللمس */\\
\quad -{}-tap-min: 44px;\\
\quad -{}-tap-comfort: 48px;\\
\quad /* المسافات على شبكة 4px */\\
\quad -{}-space-1: 0.25rem; -{}-space-2: 0.5rem; -{}-space-3: 0.75rem;\\
\quad -{}-space-4: 1rem; -{}-space-6: 1.5rem;\\
\quad /* نص متجاوب */\\
\quad -{}-fs-body: clamp(0.95rem, 0.9rem + 0.25vw, 1.05rem);\\
\quad -{}-fs-h1: clamp(1.6rem, 1.3rem + 1.5vw, 2.4rem);\\
\quad -{}-fs-h2: clamp(1.25rem, 1.1rem + 0.8vw, 1.6rem);\\
\quad /* حشو الصفحة المتجاوب */\\
\quad -{}-page-pad-x: clamp(1rem, 4vw, 2rem);\\
\quad -{}-nav-h: 64px;\\
\}
}
\end{codebox}

\begin{goalbox}
بتوحيد هذه الرموز يصبح تعديل مقياس الجوال كله من مكان واحد، وتختفي الأرقام السحرية المتفرّقة.
\end{goalbox}

% ============================================================
\section{المرحلة الثانية: توحيد مناطق اللمس}

\begin{itemize}
  \item كل عنصر تفاعلي (\texttt{button}, \texttt{a}, \texttt{input}) يحصل على حد أدنى \texttt{min-height: var(--tap-min)}.
  \item الأيقونات الصغيرة (المفضلة، المشاركة، التبويبات) تُغلّف بمنطقة لمس \texttt{44px} حتى لو كانت الأيقونة أصغر.
  \item المسافة بين عنصرين تفاعليين متجاورين لا تقل عن \texttt{--space-2}.
  \item الانتقالات تُلغى عند \texttt{prefers-reduced-motion}.
\end{itemize}

% ============================================================
\section{المرحلة الثالثة: النصوص والتخطيط المتجاوب}

\begin{itemize}
  \item استبدال \texttt{fontSize} الثابتة في الأنماط السطرية برموز النص المتجاوب.
  \item استخدام \texttt{clamp()} لعناوين الصفحات والأقسام بدل الأحجام الثابتة.
  \item \texttt{max-width} للمحتوى على الجوال، و\texttt{min()} لتحديد عرض البطاقات.
  \item ضبط \texttt{line-height} للنصوص العربية الطويلة (القرآن/الحديث) إلى $\ge 1.8$.
  \item استخدام \texttt{grid/flex} مع \texttt{gap} بدل الهوامش اليدوية.
\end{itemize}

% ============================================================
\section{المرحلة الرابعة: شريط التنقل والمنطقة السفلية}

\begin{itemize}
  \item إخفاء الشريط عند فتح لوحة المفاتيح (موجود جزئيًا) — يُعمّم على كل الحقول النصية.
  \item ارتفاع موحّد \texttt{var(--nav-h)} مع \texttt{safe-area-inset-bottom} دائمًا.
  \item المحتوى الرئيسي يحصل على \texttt{padding-bottom} كافٍ حتى لا يحجب الشريط آخر عنصر.
  \item مؤشّر الصفحة النشطة واضح وبلون الهوية.
  \item دعم الإيماءة/السحب للتبويبات (اختياري، إن ناسب).
\end{itemize}

% ============================================================
\section{المرحلة الخامسة: الأداء وإحساس التطبيق الأصلي}

\begin{itemize}
  \item \texttt{-webkit-tap-highlight-color: transparent} (موجود) و\texttt{touch-action: manipulation} للأزرار.
  \item تفادي \texttt{100vh} لصالح \texttt{100dvh} في كل الشاشات الممتدة.
  \item \texttt{content-visibility: auto} للأقسام الطويلة لتحسين التمرير.
  \item صور المسائل بتنسيق حديث (\texttt{loading="lazy"}, \texttt{decoding="async"}).
  \item ضبط لون شريط حالة iOS والثيم اللوني على هوية المنصة.
\end{itemize}

% ============================================================
\section{المرحلة السادسة: الإتاحة}

\begin{itemize}
  \item أهداف لمس \texttt{44px} كحد أدنى.
  \item دعم تكبير الخط (\texttt{--base-font-size}) الموجود + عدم كسر التخطيط عنده.
  \item \texttt{aria-label} لكل زر أيقوني بلا نص.
  \item احترام \texttt{prefers-reduced-motion} و\texttt{prefers-color-scheme}.
  \item تباين لوني كافٍ (WCAG AA) في وضع الجوال خاصة.
\end{itemize}

% ============================================================
\section{الملفات والتحقق}

\textbf{ملفات رئيسية مقترح تعديلها:}
\begin{itemize}
  \item \texttt{src/index.css} — إضافة رموز التصميم وتوحيد الاستجابة.
  \item \texttt{src/styles/exterior.css} — تقليل \texttt{!important} والتوحيد.
  \item \texttt{src/Components/Sidebar/Slider.css} — ضبط الجوال.
  \item \texttt{src/Components/UI/InstallPrompt.module.css} — مراجعة الاستجابة.
  \item المكوّنات ذات الأنماط السطرية — استبدال الأحجام الثابتة بالمتغيّرات.
\end{itemize}

\vspace{6pt}
\textbf{أوامر التحقق:}
\begin{codebox}
{\ttfamily\small\raggedleft
npm run lint\\
npx tsc -{}-noEmit\\
npm run test\\
npx playwright test   (اختبارات E2E للجوال)
}
\end{codebox}

\begin{warnbox}
\textbf{المخاطر:} تستخدم المنصة Bootstrap Class-based، وتغيير الأنماط قد يؤثّر على تخطيطات عدّة.
يُنصح بالعمل تدريجيًا صفحة-صفحة مع فحص بصري على أحجام: 360px، 390px، 430px، 768px.
\end{warnbox}

% ============================================================
\section{معايير القبول}

\begin{itemize}
  \item[${\color{okgreen}\faCheckCircle}$] كل عنصر تفاعلي لا يقلّ عن \texttt{44px} لمسًا.
  \item[${\color{okgreen}\faCheckCircle}$] لا توجد أحجام نص ثابتة تكسر التخطيط عند تكبير الخط.
  \item[${\color{okgreen}\faCheckCircle}$] التخطيط سليم على العروض: 360 / 390 / 430 / 768px.
  \item[${\color{okgreen}\faCheckCircle}$] شريط التنقل لا يحجب المحتوى ولا يتعارض مع لوحة المفاتيح.
  \item[${\color{okgreen}\faCheckCircle}$] تجربة الجوال متّسقة بين جميع الصفحات والمكوّنات.
  \item[${\color{okgreen}\faCheckCircle}$] \texttt{lint} و \texttt{tsc} واختبارات E2E تمرّ بنجاح.
\end{itemize}

\end{document}
